\begin{afterword}
\summaryhead

The post borrows the party game Taboo, in which a player must get a partner to guess a word without saying it or five related words. The author proposes the same exercise for arguments.

Two people who dispute whether an unheard falling tree makes a sound stop disagreeing once the word ``sound'' is banned: one means acoustic vibrations, the other means auditory experience. In the reverse case, two people can share a word while expecting different things, and banning the word would show it. The author suggests that futurists who agree artificial intelligence is coming, and religious believers who seem united, would find their agreement breaks down once the key words are banned.

The post concludes that in philosophical difficulties one should not define problem terms but think without them, describing observables and mechanisms. It recommends this for the free will debate, also banning words such as ``choose'' and ``responsible.'' The author calls it a nonstandard tool that works far better than defining terms.

\responsehead

The post promotes a technique and a stance together. The technique is to argue without the disputed word. The stance is the author's ``toolbox'' against what ``most philosophers would advise.'' The post claims the technique is superior, and it never makes the comparison.

The evidence is weak. The three real-world cases are futurists, believers and free will. None is carried out. The first is based on ``I would guess.'' The second announces its result, that ``mostly they won't be able to answer at all,'' before anyone is asked. For the third the author gives the instructions and does not do the exercise. The one outside source, Legg and Hutter's collection of definitions of intelligence, concludes that the definitions ``share many common features,'' which contradicts the post's use of it.

As specified here, the technique can also be used to fix the outcome. Whoever chooses the banned words can steer the result. For free will the post bans ``responsible.'' When Chalmers applied the same method in 2011, the residue he named as a sign of real disagreement was a sentence about moral responsibility. Banned in advance, it cannot appear, and the dispute will look verbal. Applied to religion, the technique comes with its verdict stated in advance. The post applies the technique only to disputes between other people and does not show it tested on a belief of the author's.

The underlying request, to state a disagreement without the disputed word, is a useful check. Chalmers later set it out together with the cases where it misleads. This post gives no such cases. Its demonstration on the syllogism repeats an objection Mill made in 1843, without credit.

In short: a method never shown working on an open question, and specified so that on the one open question it bans the words in which the real disagreement would be stated.
\end{afterword}
